\section*{8. 金属のイオンへのなりやすさ}
\begin{enumerate}
    \item 3つのビーカーA〜Cを用意し、以下の組み合わせで金属のイオンへのなりやすさを調べる実験を行った。
    \begin{itemize}
        \item[A:] 硫酸銅水溶液にマグネシウムリボンを入れると、マグネシウムの表面に赤い物質が付着した。
        \item[B:] 硫酸亜鉛水溶液に銅片を入れると、変化は見られなかった。
        \item[C:] 硫酸亜鉛水溶液にマグネシウムリボンを入れると、マグネシウムの表面に黒っぽい物質が付着した。
    \end{itemize}
    \begin{enumerate}
        \item ビーカーAで付着した赤い物質は何か。名称と化学式を答えなさい。\\[0.2cm]
        名称（ \hspace{3cm} ） \hspace{1cm} 化学式（ \hspace{3cm} ）
        \item ビーカーAの反応から、マグネシウムと銅では、どちらがイオンになりやすいといえるか。名称で答えなさい。\\[0.2cm]
        （ \hspace{5cm} ）
        \item ビーカーBの結果から、亜鉛と銅では、どちらがイオンになりやすいといえるか。名称で答えなさい。\\[0.2cm]
        （ \hspace{5cm} ）
        \item ビーカーA〜Cの実験結果をもとに、マグネシウム、亜鉛、銅を、イオンになりやすい順に左から並べ、\textbf{化学式（元素記号）}で答えなさい。\\[0.2cm]
        （ \hspace{2cm} ） $\rightarrow$ （ \hspace{2cm} ） $\rightarrow$ （ \hspace{2cm} ）
    \end{enumerate}
\end{enumerate}
